% 本文件由 LaTeX 排版 Agent 根据有效 translation.json 自动生成。
% author=Councils; work_id=3815; title=萨底迦会议（公元三​四四年）

\chapter{萨底迦会议（公元三​四四年）}

% structure: p0001 source=h1 target=omitted reason=page-title-already-rendered-by-parent

% structure: p0002 source=h2 target=section reason=relative-heading-rank; base=h2

\section{引言}

关于萨底迦会议的探讨，必须从对该会议日期的年代学考察开始。苏格拉底与索佐曼明确将其定在公元347年，并更精确地指出，该会议是在鲁菲努斯与尤西比乌斯担任执政官期间、君士坦丁大帝去世后第十一年举行的，因此按照我们的纪年方式，应在347年5月22日之后。\par

这一直是最普遍的观点，直到一百多年前，博学的西庇阿·马费伊在维罗纳发现了一份古老的亚历山大年代记（\WorkTitle{无首史}）拉丁译本的残片，并于1738年将其编辑出版在\WorkTitle{文学观察}第三卷中。该残片包含信息：在346年、君士坦提乌斯四世与君士坦斯二世担任执政官期间的法欧斐月24日（10月21日），阿塔纳修斯已从第二次流放地返回亚历山大里亚。然而，正如我们稍后将更清楚地表明的那样，人们普遍认为这次返回肯定是在萨底迦会议之后大约两年才发生的，因此曼西从中看到将这次会议日期提前至344年的必要性。圣热罗尼莫在续编的尤西比乌斯编年史中证实了这一点，他根据\WorkTitle{无首史}将圣阿塔纳修斯的返回定在君士坦提乌斯皇帝在位第十年，即346年。\par

许多学者现在追随曼西，大多数人盲目跟随；另一些人则试图反驳他，起初是博学的多明我会士马马奇，然后是韦策尔博士（弗莱堡教授），最后是我们自己在1852年\WorkTitle{蒂宾根神学季刊}上发表的一篇论文\Quote{关于萨底迦会议的争论}。不久之后又有新的发现。圣阿塔纳修斯的一些\WorkTitle{节期书信}，此前被认为已失传，在埃及的一座隐修院中被发现，附有一份非常古老的译成叙利亚语的序言，并由库雷顿在伦敦以叙利亚语出版，1852年由拉尔索教授在柏林的灰衣修士会院以德语出版。\par

在这些\WorkTitle{节期书信}中，第十九封是为347年逾越节写的，因此写于该年初，其引言明确说明是在亚历山大里亚重新抄写的。这证实了\WorkTitle{无首史}的陈述，即阿塔纳修斯已于346年10月返回亚历山大里亚，并证实了曼西假说的主要论点；另一方面，它通过阿塔纳修斯本人的证言，无可辩驳地驳斥了苏格拉底和索佐曼（由于他们相互依赖，只能算作一个证言）关于347年日期的说法。\par

如我们所说，曼西将这次会议定在344年；但圣阿塔纳修斯\WorkTitle{节期书信}的旧序言将其定在343年，事实上我们现在只能在343年和344年之间犹豫。如果序言像\WorkTitle{节期书信}本身一样古老且具有强大的说服力，那么关于萨底迦会议日期的问题就会得到最准确的确定。然而，由于该序言在几处包含错误，尤其是年代错误——例如关于君士坦丁大帝的死期——我们不能无条件地接受其关于344年的说法，只能在其与当时其他日期相符时接受。\par

无论如何，让我们假设阿塔纳修斯在340年逾越节前后到达罗马。众所周知，他在那里整整待了三年，在第四年年初被召到米兰的君士坦斯皇帝那里。这指向343年夏天。他从那里经高卢前往萨底迦，因此那次会议很有可能是在343年秋天开始的。不过，它可能一直持续到春天；因为当会议派往君士坦斯皇帝的两名特使，科隆的欧弗拉特斯和卡普亚的文森特，到达安提约基雅时，已经是344年逾越节了。该城主教斯德望以真正恶魔般的方式对待他们；但他的恶行很快昭然若揭，一次会议得以召开，在344年逾越节后将他罢免。其成员是尤西比乌斯派，因此他们任命莱昂提乌斯·卡斯特拉图斯为斯德望的继任者，这确实是阿塔纳修斯所指的那个会议，他说这次会议是在\WorkTitle{圣堂奉献会议}之后三年举行的，并制定了一份非常明确的尤西比乌斯派信仰宣认，即\OriginalTerm{长文信经}{\GreekText{μακρόστιχος}}。\par

安提约基雅主教斯德望的可耻行为一度使皇帝对亚略派减少信任，并在344年夏天允许阿塔纳修斯流亡的圣职人员回家。十个月后，伪主教亚历山大里亚的额我略去世（345年6月），君士坦提乌斯没有允许对亚历山大里亚教区进行任何新的任命，而是通过三封信召回了圣阿塔纳修斯，并等了他一年多。因此亚历山大里亚教区空置了一年多，直到346年的最后六个月。最终，在346年10月，阿塔纳修斯回到了他的主教职位。\par

那么我们看到，通过接纳圣亚大纳削的\WorkTitle{巴斯卦书信}及其序言的明确陈述，我们获得了一个令人满意的年代学体系，其中各个细节彼此吻合，因而显得合理。我们先前自己对344年日期提出的一个重大反对意见现在可以得到解决。确实，在353或354年，教宗利伯略曾写道：\Quote{八年前，欧瑟比派的代表欧多克休和玛尔提留（他们带着\OriginalTerm{长信条}{\GreekText{μακρόστικος}}信条来到西方）在米兰拒绝诅咒亚略异端。}但这里提到的米兰会议（约在345年）并非如我们先前错误假设的那样在撒底加会议之前举行，而是在其后。关于另一个困难，我们稍欠幸运。聚集在斐理伯堡的欧瑟比派（即撒底加伪会议）在他们的会议信函中说：\Quote{迦萨主教阿斯克勒帕斯在十七年前被免去主教职务。}这次免职发生在一次安提约基雅会议上。如果我们将此会议等同于著名的330年会议（安提约基雅的欧斯塔提也在那次会议上被推翻），那么往后推算十七年，就会得到346或347年，即斐理伯堡会议信函的写作年份，因此也是撒底加会议的年份。然而，有两种方式可以避免这一结论：要么我们必须假定阿斯克勒帕斯在330年安提约基雅会议前一年左右已被免职，要么认为斐理伯堡会议信函的拉丁译本（我们不再拥有原文）中关于数字十七的陈述是笔误或手误。但无论如何，这封会议信函无法改变这样一个事实：亚大纳削在写作347年巴斯卦书信时已回到亚历山大里亚，因此撒底加会议必定是在此之前几年举行的。\par

% structure: p0011 source=h3 target=subsection reason=relative-heading-rank; base=h2

\subsection{关于教规文本的说明}

撒底加教规的希腊文本和拉丁文本均已流传下来；一些作者如\Person{里歇尔}（\WorkTitle{大公会议史}，卷一，第98页）认为拉丁文本才是原文，而其他作者如\Person{瓦尔希}（\WorkTitle{教会会议史}，第179页）则得出了完全相反的结论。然而，现在主要归功于\Person{巴莱里尼}和\Person{斯皮特勒}的研究，学者们的一致意见——如\Person{赫费勒}所言——是这些教规最初以两种语言起草，因为它们是为拉丁人和希腊人共同制定的。或许我可以提醒读者，在许多西方教规汇编中，撒底加教规紧接在尼西亞教规之后，没有任何间断，也没有注明它们并非在该会议上制定。同样值得牢记的是，这些教规在特鲁洛会议上被希腊人视为具有普世会议的权威，因此被纳入希腊教会法的主体中。\par

我已为读者提供了每种文本的非常准确的翻译。\par

% structure: p0014 source=h2 target=section reason=relative-heading-rank; base=h2

\section{教规}

从各省聚集在撒底加的圣会议规定如下。\par

（见于君士坦丁堡的若望六世纪汇编及其他若干手稿的希腊文本。亦见于希腊注释家的著作。拉丁文本见于\WorkTitle{古代本}、\WorkTitle{狄奥尼削·爱西古译本}以及\WorkTitle{伊西多尔译本}，包括真本和伪本。）\par

% structure: p0017 source=h3 target=subsection reason=relative-heading-rank; base=h2

\subsection{第一条}

科尔多瓦城主教奥修斯说：必须从根基上铲除一种普遍的恶行，或者说一种极为有害的腐败。绝不允许任何主教从小城迁移到另一城，因为这种过失的明显原因可以解释这类企图：从未发现任何主教试图从大城迁到小城。因此，显然此类人受到过度贪婪的驱使，只为追求权势之名而服务野心。那么，如此严重的滥用是否应受到严厉惩罚？因为我认为，这种人甚至不应被接纳进入平信徒共融。全体主教说：这是共同的意愿。\par

\Person{奥修斯}主教说：必须从根基上铲除一种普遍的恶行和有害的腐败。绝不允许任何主教从自己的城迁移到另一城。因为此类企图的理由是明显的，因为在此事上，从未发现任何主教从大城迁到小城的。因此，显然这些人被过度的贪婪所激动，是在服务野心，并且旨在获取权力。如果这是共同的意愿，就应极其严厉和严格地惩罚如此大的恶行，以致这样的人甚至不被接纳进入平信徒共融。全体一致回答：这是我们的意愿。\par

% structure: p0020 source=h3 target=subsection reason=relative-heading-rank; base=h2

\subsection{第二条}

主教奥修斯说：但如果发现任何人如此疯狂或鲁莽，以致企图以上诉为由声称他从会众（平信徒）带来了信件，那么显然，少数被贿赂和报酬败坏的人可能会在教会中制造骚动，要求，诚然，要求该人为主教。我想，这种伎俩和设计必须受到惩罚，以致这样的人甚至被认为不配得到临终平信徒共融。因此，请回答这一判决是否得到你们的认可。他们（主教们）回答：所说的话已获认可。\par

\Person{奥修斯}主教说：即使任何人如此鲁莽，以致或许以上诉为由声称他从会众收到了信件，因为显然，少数人可能被报酬和贿赂败坏——即那些不持守纯正信仰的人——在教会中制造骚动，并似乎要求该人为主教；我判定这些欺诈必须受到谴责，以致这样的人在临终时也不应得到平信徒共融。如果你们都同意，就颁布它。会议回答：我们同意。\par

\Person{赫费勒}。\par

在拉丁文译本中，加上了\Emph{\OriginalTerm{那些没有真诚信仰的人}{qui sinceram fidem non habent}}，这一短语见于\Person{狄奥尼修斯·埃克西古斯}、\Person{伊西多尔}以及《旧拉丁译本》（\OriginalTerm{《普里斯卡》}{Prisca}）中。其含义如下：\Quote{在某个城镇中，少数人——尤其是那些没有真信仰的人——容易被贿赂，要求某人作主教。}萨迪卡会议的教父们显然在此处指向\OriginalTerm{亚略派}{Arians}及其追随者，他们通过这类阴谋，一旦在一个城镇中赢得了哪怕一小派的支持，就试图强行进入主教职位。此外，\Place{安提约基雅}会议在公元341年，虽然严格来说\OriginalTerm{欧瑟比派}{Eusebians}在其中占主导地位，却在其第二十一号法令中制定了一条类似但较为宽松的规则。\par

% structure: p0025 source=h3 target=subsection reason=relative-heading-rank; base=h2

\subsection{第三条}

\Person{何西乌}主教说：还必须补充这一条——任何主教都不得从自己的省份前往已有主教的另一个省份，除非他被他的弟兄们邀请，免得我们似乎关闭了爱德的大门。\par

同样也要预见到这种情况：如果在任何省份中，一位主教对他的弟兄和同为主教的人有某种事务，他们两人中的任何一方都不应请来另一省份的主教作为仲裁人。\par

但如果遇有人在任何事务上被判有罪，而他认为自己的案件并非无理而是合理，为使这个问题得以重新审理，那么，若你们爱德以为妥当，让我们敬礼宗徒伯多禄的纪念，并让作出判决的人写信给罗马主教儒略，以便必要时该案件可由邻近省份的主教们复审，并由他指定仲裁人；但如果不能证明他的案件属于需要重新审理的性质，则一经作出的判决不应废止，而应如前一样保持有效。\par

\Person{何西乌}主教说：还必须补充这一条——任何主教都不得从自己的省份前往已有主教的另一个省份，除非偶然受到他们的弟兄邀请，免得我们似乎关闭了爱德之门。\par

但如果在任何省份中，一位主教与他的弟兄主教有争议，两人中的任何一方都不应请来另一省份的主教作为裁判者。\par

但如果遇有人在任何案件中被判有罪，而他认为自己有合理理由，为使问题得以重新审理，那么，若你们同意，让我们敬礼宗徒圣伯多禄的纪念，并让审理案件的人写信给罗马主教儒略；如果他认为该案件应当重审，就这样做，并由他指定裁判者；但如果他发现该案件的性质无需推翻原先的判决，则他所决定的事应被确认。这是否为各位所同意？会议回答：我们同意。\par

% structure: p0032 source=h3 target=subsection reason=relative-heading-rank; base=h2

\subsection{第四条}

\Person{高登蒂}主教说：如果你们认为妥当，有必要在你们所宣布的这充满真诚爱德的决议上补充一点：如果任何主教被这些邻近主教的判决所撤职，并声称他有新的辩护理由，则不得在罗马主教就此作出判断和裁决之前，在其教座上另立新主教。\par

\Person{高登蒂}主教说：如果你们乐意，有必要在你们所宣布的这充满圣善的判决上补充一点——当任何主教已被那些在邻近地方拥有教座的主教的判决撤职，而他（被撤职的主教）宣称他的案件将在罗马城中审理时——则在该显然已被撤职的人提出上诉之后，决不可另立另一主教为其教座，除非该案件已由罗马主教的判决所确定。\par

% structure: p0035 source=h3 target=subsection reason=relative-heading-rank; base=h2

\subsection{第五条}

\Person{何西乌}主教说：法令规定，如果任何主教被控告，同一区域的主教们聚集并撤销其职务，而他（近似说）上诉并投靠罗马教会最蒙福的主教，并且该主教愿意听取他的申诉，认为应当重新审查他的案件，那么他应乐意写信给那些最靠近该省份的、同为主教的人，以便他们仔细而准确地审查细节，并根据真理之言对这件事进行投票。如果任何人要求再次审理他的案件，并且在他的请求下，罗马主教认为适宜派遣\OriginalTerm{来自他身旁的}{a latere}司铎，那么该主教有权根据他认为良好和正当的判断——派遣一些人，与主教们一起担任审判者，并授予他们派遣者的权威。这一点也应加以规定。但如果他认为主教们足以审理和决定此事，则他可根据自己最明智的判断行事。\par

众主教回答：所说的话获得批准。\par

\Person{何西乌}主教说：此外，法令规定，如果一位主教被控告，该区域的主教们聚集并撤销其职务，如果被撤职的人上诉并投靠罗马教会的主教，并希望得到审理，若他认为应当重新审理或检查他的案件，他应乐意写信给邻近省份的主教们，让他们仔细查问所有细节，并依据真理之言作出决定。但如果请求重审案件的人通过恳求促使罗马主教派遣一位\OriginalTerm{来自他身旁的}{a latere}司铎，则该主教有权决定他所要执行和确定的事；如果他决定派遣一些人，他们应到场并与主教们一同担任审判者，且被授予派遣者的权威，则应按他的选择行事。但如果他认为主教们足以作出最终决定，则他可根据自己最明智的判断行事。\par

% structure: p0039 source=h3 target=subsection reason=relative-heading-rank; base=h2

\subsection{第六条}

奥西乌斯主教说：\Quote{如果一个教省中主教人数众多，而某一位主教缺席，或因疏忽未前来参加主教会议，且对主教们所作出的任命没有表示同意，但民众集会恳求为他们所期望的主教举行祝圣——则有必要先由该教省的督主教（我指的自然是大都会的主教）致函提醒那位缺席的主教，说明民众要求为他们派遣一位牧者。我认为，等待那位缺席主教到来也是妥当的。但如果经信函召唤后他仍未前来，甚至没有书面答复，则应当顺从民众的意愿。}\par

在祝圣都主教时，也应邀请邻近教省的主教们前来参加。\par

绝对不许在村庄或小镇上祝圣主教，因为即使只有一位司铎也绰绰有余（在那里并无必要祝圣主教），以免主教的名衔和权威遭受轻慢。但如前所述，教省的主教们应在先前已有主教的那些城市中祝圣主教；若发现某城市人口众多，足以配得上设立主教座，则可给予一位主教。\par

\Quote{各位是否同意？}众人回答：\Quote{我们同意。}\par

奥西乌斯主教说：\Quote{如果在一个教省中原本有很多主教，如今只剩一位主教，而他因疏忽未选择祝圣主教，且民众已提出请求，则邻近教省的主教们应首先致函居住在该教省的那位主教，表明民众为自己寻求一位治理者（即牧者），并指出这是合理的，以致他们也前来与他一同祝圣主教。但如果他拒绝承认他们的书面通讯，置之不理，且不作书面答复，则应当满足民众的请求，以便邻近教省的主教们前来祝圣一位主教。}\par

但不允许在任何一个村庄或不重要的城市中祝圣主教，因为一位司铎便已足够，以免主教的名衔和权威变得廉价。那些从另一教省被邀请来的主教们，除非在那些先前已有主教的城中，或在如此重要或人口众多以至于有权拥有主教的城中，否则不应祝圣主教。\par

\Quote{各位是否同意？}主教会议回答：\Quote{我们同意。}\par

本条法规的第二部分在拉丁文中完全缺失。希腊评注家佐纳拉斯、巴尔萨蒙和阿里斯泰努斯理解此规定是指\Quote{在任命都主教时，也应邀请邻近教省的主教们前来，\InnerQuote{大概是为了给这一行为增添更多的庄严，}}赫费勒如是说。范埃斯彭、蒂勒蒙和赫尔布斯特也同意此说。\par

希腊文本与拉丁文本在第一部分含义不同：希腊文本考虑的是在选举填补空缺时，一位主教从主教会议中缺席的情况；拉丁文本考虑的则是一个教省中（在战争、瘟疫或类似情况后）仅剩一位主教的情况。这第二种含义被范埃斯彭、克里斯蒂安·卢普斯等人接受。此外，从弗洛多亚德的\WorkTitle{兰斯教会史}（\WorkTitle{兰斯教会史}，第三卷，第二十章〔我从未见过的书〕）来看，高卢教会似乎是依据对本条法规的这一理解来行事的。格拉提安也持此见。\par

在拉丁文本与希腊文本之间，还有佐纳拉斯的解释，即如果一个教省此前曾有许多主教，但后来因某种变故仅剩一位主教（都主教除外），而他又疏忽出席新主教的祝圣，则应由都主教函召他前来；若他届时未到，则祝圣可不受其阻碍地进行。哈尔梅诺普洛斯也同意此解释，并进一步指出都主教可独自祝圣主教，他依据\OriginalTerm{}{\GreekText{τὸ} \GreekText{ἱκαυὸν} \GreekText{κ}}. \OriginalTerm{}{\GreekText{τ}} \OriginalTerm{}{\GreekText{λ}}. 这段话。\par

有些学者认为现行的希腊文本和拉丁文本均非原版，但希腊文本较为接近原版，不过须根据马费伊在维罗纳的一份古抄本中发现的一种古老拉丁译本予以修正。巴莱里尼兄弟在其对\WorkTitle{圣大良文集}的注释中（第三卷第xxxij页第4节）曾仔细关注这一点。看来，382年君士坦丁堡的教父们所引用的可能就是本条法规；如果是这样，那么当时他们所用的希腊文本似乎与维罗纳译本所据的希腊文本相同。\par

% structure: p0051 source=h3 target=subsection reason=relative-heading-rank; base=h2

\subsection{第七条}

奥西乌斯主教说：\Quote{我们的纠缠不休、极其固执以及不公正的请求，导致我们失去了应有的恩惠和信任。因为许多主教不断前往朝廷，尤其是非洲的主教们——正如我们从我们亲爱的兄弟和同为主教的格拉图斯所获悉的——他们不接受有益的劝告，反而轻视它们，以至于有人将众多且五花八门、对教会毫无益处的请愿书带到朝廷，并没有（如应当且适宜地）帮助和支持穷人、平信徒或寡妇，反而图谋为某些人谋取世俗的尊荣和职位。这种恶行因此引发了怨言（\OriginalTerm{更好的是，怨言}{\GreekText{τονθρυσμόν}}或\OriginalTerm{}{\GreekText{τονθορυσμόν}}），同时伴随着对我们的某种丑闻和指责。但我认为，主教帮助受压迫者、遭受不公的寡妇、或被剥夺遗产的孤儿是完全恰当的，只要这些人有正当的申诉理由。}\par

那麼，親愛的弟兄們，若此舉為眾人所贊同，便規定：除非我們至虔誠的皇帝以其親筆書信所傳召者，不得有主教前往朝廷。然而，由於常有因其過犯被判處流放或放逐海島、或受其他判決之人，懇求寬恕並逃往教會（即避難所），此等人不應被拒絕援助，而應毫不遲延、毫不猶豫地為他們請求赦免。若此舉亦為你們所同意，請全體表決贊同。\par

眾人回答說：也請如此規定。\par

主教\Person{奧西烏斯}說：強求、過度的固執及不正當的請求，使我們甚少獲得恩寵或信任，而某些主教卻不斷前往朝廷，尤其非洲人，如我們所知，他們輕視並蔑視我們至聖的兄弟及同僚主教\Person{格拉圖斯}的救世勸告，以致他們不僅向朝廷呈遞許多不同的請求（並非為教會的利益，亦非如常理那般，為救助窮人、寡婦或孤兒），甚至為某些人謀求世俗的尊榮和官職。因此，此惡行時常不僅引發怨言，甚至引起惡表。然而，主教理應為遭受暴力與壓迫、受困苦的寡婦及被欺詐的孤兒代禱，惟其前提是這些人有正當的理由或請求。\par

那麼，親愛極愛的弟兄們，若你們如此贊同，我們規定：除非受敬畏天主的皇帝書信邀請或傳召者，不得有主教前往朝廷。然而，由於常有遭受不公義者，或因其過犯被判處流放或放逐海島，或總之受其他判決之人，逃往教會的慈懷，此等人應得到救助，並應毫不猶豫地為他們請求赦免。請你們贊同此規定。\par

眾人說：我們贊同，請即規定。\par

% structure: p0058 source=h3 target=subsection reason=relative-heading-rank; base=h2

\subsection{第八條}

主教\Person{奧西烏斯}說：也請你們的明智決定此事——鑑於為使主教不致因前往朝廷而受責備而如此規定——若有人具有我們上述所說的請求書，應由其一位執事代為呈遞。因為下屬的身份不會引起嫉妒，而所獲准之事（由皇帝）便能更快地回報。\par

眾人回答說：也請如此規定。\par

主教\Person{奧西烏斯}說：也請你們的遠見如此規定——鑑於你們已作此規定，以免主教之膽大妄為（或：被看到）前往朝廷。因此，凡有或收到如上所述之請求書者，應由其一位執事代為呈遞，因為事奉者的身份不會是嫉妒的對象，而他將能更快地回報所獲得之事。\par

\Signature{Van Espen}\par

此規定有三重含義。第一，主教前往朝廷時，不應在朝廷或在自己人民中引起懷疑，以為他接近君主是為獲取自己之事。第二，根據\Person{佐納拉斯}的解釋：\Quote{無人應對留在軍營中的事奉者或執事生氣，因主教已離開那裡}。第三，事奉者能帶走他所求之事，即（按\Person{佐納拉斯}所言）皇帝赦免過錯的書信，或類似其他事項。\par

% structure: p0064 source=h3 target=subsection reason=relative-heading-rank; base=h2

\subsection{第九條}

主教\Person{奧西烏斯}說：我認為，此舉亦隨之而來：若在任何省份內，主教們向其兄弟及同僚主教之一呈送請求書，則居於最大城市（即府城）者，應親自派遣其執事及請求書，並為其提供推薦書信，當然也依次致信給我們的兄弟及同僚主教，若其時有人正在最虔誠的皇帝管理公務的地方或城市停留。\par

但若某位主教在朝廷有朋友，且欲就某正當目的向他們請求，不應禁止他透過其執事提出請求，並促使這些朋友按其意願給予仁慈協助。\par

而那些來到\Place{羅馬}的人，如我之前所說，應將他們所要呈遞的請求書交付給我們親愛的兄弟及同僚主教\Person{尤利烏斯}，以便他先加審查，以免其中有些請求不當，然後憑藉他的倡導和關懷，將請求書送往朝廷。\par

全體主教回答說：這是他們的意願，且此規章最為恰當。\par

此舉也隨之而來：從任何省份，若主教們向那位在府城有主教座位的我們兄弟及同僚主教呈送請求書，他（府主教）應派遣其執事攜帶請求書，並為其提供內容相仿的推薦書信，致給其時居住在那位有福且受祝福的皇帝治理國家之地區與城市的我們的兄弟及同僚主教。\par

然而，若某位主教欲求取某項請求（即正當者）而在宮中有朋友，則不禁止他透過其執事提出請求，並告知那些他知道能在其缺席時仁慈為他代禱的人。\par

十、但那些來到\Place{羅馬}的人，如前所述，應將他們所攜帶的請求書交付給我們至聖的兄弟及同僚主教、羅馬教會的主教，以便他也審查這些請求是否正當公義，並付出勤勉與關懷，使它們被轉送至朝廷。\par

眾人說：這是他們的意願，且此規章恰當。\par

主教\Person{阿利皮烏斯}說：若他們為孤兒、寡婦或任何困苦者且擁有並非不公義的案件之故，而承受旅途勞頓，他們將有正當理由（為其行程）；但如今他們主要提出那些無法不引起嫉妒和責備的請求，則他們無需前往朝廷。\par

此處的拉丁文本不僅是希臘文本的翻譯，更是一種詮釋，因為它明確指出每位主教應將他打算呈給朝廷的請求書首先送給他的府主教，由後者遞交。這在希臘文本中並不明確，然而希臘注釋家卻在其中找到此意。\par

赫费勒（\Person{Hefele}）：\Quote{佐纳拉斯、巴尔萨蒙与阿里斯泰诺斯对此法令解释略有不同，即：\InnerQuote{若一位主教希望将其上呈皇帝的请愿书送至皇帝驻跸城市的主教，他应先将请愿书呈交该省的总主教（据阿里斯泰诺斯，即其本省总主教），然后由总主教派遣其自己的执事，持推荐函致可能在朝廷的主教（们）。}此分歧源于对法令开篇的\OriginalTerm{弟兄与同僚主教}{brother and fellow bishop}的不同理解。我们以此指其本省总主教，并将\OriginalTerm{在更大的……等等}{\GreekText{ὁ} \GreekText{ἐν} \GreekText{τῇ} \GreekText{μείζονι} \GreekText{κ}.\GreekText{τ}.\GreekText{λ}.}视为对\OriginalTerm{同僚主教}{fellow-bishop}的更精确界定，将分词\OriginalTerm{\GreekText{τυγχάνων}}{\GreekText{τυγχάνων}}视为与\OriginalTerm{\GreekText{τυγχάνει}}{\GreekText{τυγχάνει}}等价，并从\OriginalTerm{他亲自连同执事}{\GreekText{αὐτὸς} \GreekText{καὶ} \GreekText{τὸν} \GreekText{διάκονον}}开始主句。贝弗里奇以同样方式翻译了此法令。相反，佐纳拉斯等人则将\OriginalTerm{同僚主教}{fellow-bishop}理解为当时皇帝驻跸地的主教，并将\OriginalTerm{在更大的……等等}{\GreekText{ὁ} \GreekText{ἐν} \GreekText{τῇ} \GreekText{μείζοη} \GreekText{κ}.\GreekText{τ}.\GreekText{λ}.}视为并非对前述内容的澄清，而是主句，意即\InnerQuote{于是总主教要……}等。根据此解释，法令中完全缺少主教必须依靠总主教的含义。}\par

此法令的第一部分为拉丁文本第九条法令的最后部分。最后部分为拉丁文本第十条法令，但关于\Emph{阿利皮乌斯}的个人部分在希腊文本中被省略。\par

% structure: p0077 source=h3 target=subsection reason=relative-heading-rank; base=h2

\subsection{第十条法令}

奥西乌主教说：我也认为这是必要的。你们应当以充分的周全和谨慎考虑：若有某富贵之人或专业律师被推举为主教，则在他履行完读经员、执事与长老之职之前，不可被祝圣；如此，他经过逐级晋升后（若堪当），方可达到主教的高位。他在每一等级中必须停留足够长的时间，以便其信德、名誉生活、品格的坚毅及举止的持重为人所知，且他、堪称享有神圣司祭职（即主教职）者，方可获得最高荣誉。因为，仓促或轻率地行事，以便急促地祝圣主教、长老或执事，是不适宜的，也不合纪律或良好操行；如此，他会被正当地视为初学者，尤其是最可福的宗徒，外邦人的导师，显然禁止了急促的祝圣；因为即使是极长时间的考验，也并非不合理地要求用以展示每位候选人的品行和性格。\par

全体一致表示，这是他们的意愿，并且此事必须绝对不可推翻。\par

奥西乌主教说：我也认为你们必须最仔细地考虑这一点：若偶尔有某富贵之人或专业律师或退休官员被推举为主教，则在他履行完读经员之职及执事与长老之职前，不可被祝圣；如此，他若显示自己堪当，便逐级上升至主教高位。因为通过这些需要相当长时间的晋升，可以考验他的信德、明辨、庄重与谦逊。若他被发现堪当，便应享有神圣的司祭职（即主教职）。因为，急促或轻率地祝圣一位初学者为主教、长老或执事，是不适宜的，也不合次序或纪律，尤其是可福的宗徒，外邦人的导师，显然明确禁止了这一点。那些其生活经过考验、其功绩经时间证明的人，才应被祝圣。\par

全体一致表示，这是他们的意愿。\par

% structure: p0082 source=h3 target=subsection reason=relative-heading-rank; base=h2

\subsection{第十一条法令}

奥西乌主教说：我们还应当规定：当一位主教从一城到另一城，或从一省到另一省，为满足虚荣，服务于自我赞美而非敬虔，并希望延长停留时，而该城的主教不善于教导，则这位来访主教不应藐视当地主教，试图通过频繁的讲道贬低他、减少他的声望（因为这种做法常引起骚乱），并以此等手段谋求并夺取他人的主教座，不顾忌放弃托付给他的教会而请求调任至另一教会。在此类情况下，应设定一个明确的时间限制，尤其因为不接受来访主教常被视为粗鲁无礼。你们还记得，先辈们曾规定：若一位平信徒在一城停留，连续三个主日（即整整三周）不参加神圣敬礼，应将他逐出共融。既然对平信徒已有此规定，那么一位主教，除非有重大需要或困难事务，长时间离开自己的教会并令托付给他的人民受苦，既非必要、不适宜，也甚至不合宜。\par

全体主教说：我们决定此规定也最为恰当。\par

主教\Person{奥西乌斯}说：你们也应当决定这一点。如果一个主教从一座城到另一座城，或从他的省份到另一个省份，出于野心而非虔诚，想要长期逗留在陌生的城中，然后（因为恰巧当地的主教不如他那样有经验或有学问），他这个外人就开始侮慢当地主教，多次发表讲话贬低他、减少他的声誉，毫不犹豫地用这种手段离开所托付给他的教会而迁移到别人的教会——那么你们\Quote{在此情形下}应为\Quote{他在城中的停留}设定一个时间限制，因为一方面拒绝接待一位主教是不礼貌的，另一方面他停留太久是有害的。必须对此作出规定。我记得在一次以前的会议上，我们的弟兄们规定，如果任何平信徒在他所停留的城中连续三个主日（即三周）不参加神圣礼仪，他就应被剥夺共融。如果对平信徒已有此规定，那么对于一个主教来说，如果没有重大必要留住他，他离开自己的教会超过上述时间就更不合法和不合适了。\par

全体说：这是他们所乐意的。\par

% structure: p0087 source=h3 target=subsection reason=relative-heading-rank; base=h2

\subsection{第十二条}

主教\Person{奥西乌斯}说：既然任何情况都不应被忽略，也应当规定这一点。我们的一些弟兄和同为主教者，在他们的主教辖区城市中财产甚少，但在其他地方拥有大量产业，而且他们还能以此帮助穷人。我想应当允许他们，如果他们要巡视自己的庄园并关注收获，可以在自己的庄园上度过三个主日（即三周），并在最近的有司铎举行礼仪的教堂参与神圣敬拜并举行礼仪，这样他们就不显得缺席了敬拜，也不致过于频繁地来到有主教的城中。这样，他们的私事不会因缺席而受损，他们也不会被指责为野心和傲慢。\par

所有主教说：这项法令也为我们所认可。\par

主教\Person{奥西乌斯}说：既然任何情况都不应被忽略\Quote{也应当规定这一点}。我们有些兄弟主教，并不居住在任命他们为教区的城中，或因为在那里的财产很少，而在别处有大量产业，或可能出于对亲族的爱和对他们的迁就。应当允许他们，去自己的庄园监督和处理收获，并且\Quote{为此目的}如果必要，可以在自己的庄园上度过三个主日（即三周）；或者，如果邻近有有司铎的城市，为了不显得在无教堂的情况下度过主日，他们可以到那里去。这样，他们的私事不会因缺席而受损，他们也不致因频繁往有主教常驻的城中而招致野心和谋职的嫌疑。全体说：这为他们所认可。\par

% structure: p0091 source=h3 target=subsection reason=relative-heading-rank; base=h2

\subsection{第十三条}

主教\Person{奥西乌斯}说：愿这也为全体所乐。如果任何执事或司铎或任何圣职人员被绝罚，并投奔另一位知道他并知道他被自己的主教开除共融的主教，\Quote{那位主教}不可因接纳他共融而得罪其兄弟主教。如果有人胆敢如此做，让他知道，他必须在一主教会议前自呈并作出交代。\par

所有主教说：这项决定将始终确保和平，并维护全体的和谐。\par

主教\Person{奥西乌斯}说：愿这也为全体所乐。如果任何执事或司铎或任何圣职人员被自己的主教拒绝共融，并去到另一位主教那里，而他所投靠的那位主教知道他已被自己的主教拒绝，那么他不可准予他共融。但如果他这样做，让他知道，他必须在一主教会议前作出交代。\par

全体说：这项决定将维护和平，保持和谐。\par

% structure: p0096 source=h3 target=subsection reason=relative-heading-rank; base=h2

\subsection{第十四条}

主教\Person{奥西乌斯}说：我必须不可不提到一个常常催迫我的问题。如果发现一个主教易于发怒（这样的人不该有此情绪），并且他突然对一位司铎或执事动怒，想要把他逐出教会，必须规定，这样的人不可被过于仓促地定罪（或读\OriginalTerm{无辜者}{\GreekText{ἀθῶον}}，如果他是无辜的）以及被剥夺共融。\par

全体说：被逐出者应被授权投奔同省都主教。如果都主教不在，他应前往最近的主教那里，请求仔细审查他的案件。因为请求审理的人不应被拒绝。\par

那位将某人逐出的主教，无论是否公正，都不应因案件被审查及其决定被确认或修改而感到不悦。但是，在所有细节被仔细忠实地审查之前，被排除共融的人不应在案件判决之前要求共融。如果任何与会的圣职人员清楚看出他身上的傲慢和自大，既然不应容忍无礼或不义的非议，他们应当用相当严厉和痛心的话来纠正这样的人，使人们顺从并服从恰当和正确的命令。因为主教应当向属下显示真诚的爱和关怀，同样，属下也应以真诚向他们的主教履行其职务的职责。\par

主教\Person{奥西乌斯}说：我必须不可不提到一个进一步触动我的问题。如果某个主教偶然易于发怒（这不应是这种情况），并急怒之下对某位司铎或执事，想要把他逐出教会，必须规定，一个无辜的人不应被定罪或剥夺共融。\par

因此，被逐出者应被授权向邻近的主教们上诉，让他的案件被更勤勉地聆听和审查。因为请求审理的人不应被拒绝。\par

那将他逐出的主教，无论有理无理，都当耐心忍受此事得以讨论，以使他的判决或得到众多【审判官】的认可，或予以修正。然而，在一切细节得到慎重而忠实的审查之前，任何人都不应在案件判决之前，擅自接纳那被排除共融者。但若参与审理者在【此类】圣职人员身上观察到傲慢与骄矜——因为主教受屈或受辱确实不合宜——就应用严厉的言辞纠正他们，使他们服从主教正当与正确的命令。因主教当向圣职人员表现真诚的慈爱与爱德，同样，其辅理人员也当向主教表现真诚的服从。\par

这是拉丁版第十七条教规。\par

% structure: p0104 source=h3 target=subsection reason=relative-heading-rank; base=h2

\subsection{教规第十八条（拉丁文）}

主教\Person{雅努阿里乌}说：愿你们的圣德也规定此事：任何主教都不得试图为他自己获取另一城市主教的教会中的辅理人员，并祝圣他到自己辖下的某一堂区。\par

众人说：我们赞成此事，因为这事上的争执容易产生不和，因此我们众人均禁止任何人胆敢这样做。\par

% structure: p0107 source=h3 target=subsection reason=relative-heading-rank; base=h2

\subsection{教规第十五条}

主教\Person{奥西乌}说：我们也当全体规定：若任何主教未经其主教同意，从另一教区祝圣他人的辅理人员至任何等级，该祝圣应视为无效且不予确认。若有人擅自如此行，应由我们的弟兄及同僚主教予以告诫和纠正。\par

众人说：这一规定也应保持不变。\par

主教\Person{奥西乌}说：我们全体也规定：若任何【主教】未经其主教同意和意愿，从另一教区祝圣他人的辅理人员，其祝圣不得被批准。凡擅自如此行的人，应由我们的弟兄及同僚主教予以告诫和纠正。\par

这是拉丁版第十九条教规。\par

% structure: p0112 source=h3 target=subsection reason=relative-heading-rank; base=h2

\subsection{教规第十六条}

主教\Person{阿埃提乌}说：你们并非不知\Place{得撒洛尼}这座都会城市何其重要何其广大。因此，司铎与执事常从其他省份来到该城，不满足于短期停留，而是留下并以此为永久居所，或者经长时期耽搁后方才勉强被迫返回自己的教会。对此事应制定一项法令。\par

主教\Person{奥西乌}说：那些为主教所制定的规定，也当为这些人遵守。\par

主教\Person{阿埃提乌}说：你们并非不知\Place{得撒洛尼}城何其广大重要。司铎与执事常从其他地区来到该城，不满足于短期停留，而是要么定居于此，要么至少经过长久拖延后方才勉强被迫返回自己的地方。\par

众人说：那些为主教所规定的期限，也当为这些人遵守。\par

% structure: p0117 source=h3 target=subsection reason=relative-heading-rank; base=h2

\subsection{教规第十七条}

此外，因我们弟兄\Person{奥林匹乌}的建议，我们也乐意规定此事：若一位主教遭受暴力并被不义驱逐，无论是因为其纪律，或是因为他对【天主教会的信仰】\OriginalTerm{宣认}{confession}，或是因为他维护真理，且虽清白虔诚却逃离危险，来到另一城市，则不应禁止他留在那里，直到他被恢复原职，或直至他能从所遭受的暴力与不义中得到解救。因为受不义驱逐者若不被我们接纳，实属残酷且极难忍受；这样的人理当受到最大的关怀与热诚接待。*\par

众人说：我们也赞成此事。\par

因我们弟兄\Person{奥林匹乌}的建议，我们也乐意规定此事：若有人遭受暴力，并因纪律及其天主教宣认，或因维护真理而不义地被驱逐，且虽清白虔诚却逃离危险，来到另一城市，则不应禁止他留在那里，直到他能返回或其冤屈得到昭雪。因拒绝接待受迫害者，实属苛刻无情；实应向他表现极大的热诚与丰厚的关怀。\par

全体主教会议说：所有已经规定之事，普世天主教会将予以保存与持守。\par

而来自各省的全体主教聚集后签名如下：\par

我，N.，\Place{N.}城及\Place{N.}省的主教，如上所写，如此相信。\par

这是拉丁版第二十一条教规，也是最后一条。\par

% structure: p0125 source=h3 target=subsection reason=relative-heading-rank; base=h2

\subsection{教规第十八条（希腊文）}

主教\Person{高登提乌}说：\Person{阿埃提乌}兄弟，你知道自从你被立为主教以来，和平一直统治【你的教区】。为使关于圣职人员的纷争不留痕迹，似乎应接纳所有由\Person{穆赛乌}和\Person{尤提基阿努}所祝圣者，只要在他们身上找不到过错。\par

% structure: p0127 source=h3 target=subsection reason=relative-heading-rank; base=h2

\subsection{教规第十九条（希腊文）}

主教\Person{奥西乌}说：以我卑微【即不配】的意见判决如下——因为我们应当对众人温和、忍耐并恒常怀有怜悯——凡曾被我们某些弟兄晋升至教会圣职的人，若他们不愿回到被委派【或：聘定】的教会，则将来不再接纳；且\Person{尤提基阿努}不应继续自称为主教，\Person{穆赛乌}也不得被视为主教；但若他们寻求平信徒共融，则不应拒绝。\par

众人说：我们赞成。\par

显然，这两条教规不存在于拉丁文本中，是因为它们不适用于拉丁教会，且仅包含针对\Place{得撒洛尼}的特殊规定。\par

% structure: p0131 source=h3 target=subsection reason=relative-heading-rank; base=h2

\subsection{教规第二十条}

高登提乌主教说：这些被我们\OriginalTerm{主教}{\GreekText{τῶν} \GreekText{ἱερέων}}视为合宜、正当、适切地制定的法令，既悦乐天主又悦乐人，若不加[惩罚的]恐惧，便不能获得应有的效力和有效性。因为我们自己知道，由于少数人的无耻，\OriginalTerm{主教}{\GreekText{τῆς} \GreekText{ἱερωσύνης}}这一神圣可敬的称号常常受到谴责。所以，如果有人出于傲慢和野心而非寻求悦乐天主，胆敢违反全体议决而采取不同行动，他应当预先知道，他必须就此事答辩和辩护，并失去主教之尊荣和职衔。\par

众人回答：此判词恰当且正确，我们亦如此乐意。\par

此法令若要广泛知晓并得以最佳实施，则凡我们当中以通衢大道为教座的主教，若看见一位主教经过，应当询问他行程的原因和目的地。若在他离开时，发现他是前往宫廷，则当按前述[前往宫廷的]目的进行查询。若他是应召而来，则不应阻碍其离开。但若他是出于炫耀（如贵爱德前述）试图前往宫廷，或为某些人请愿，则既不得签署他的信函，也不得收纳他进入共融。\par

众人都说：也当如此定断。\par

高登提乌主教说：你们在[法令中]合宜且适当地制定的这些事项，在众人（或：向众人）看来，既悦乐天主（或：又悦乐人），唯有在此行动中加入[惩罚的]恐惧时，才能获得效力和有效性。因为我们自己知道，由于少数人的无耻，神圣可敬的司祭（即主教）之名曾多次受到责备。所以，若有人胆敢反对全体之判决，且寻求服侍野心而非悦乐天主，则必须让他知道，他将要答辩并失去职位和品级。\par

唯有当我们中那些教座位于大道上的各位，若看见一位主教经过，询问他旅程的原因，查明他的目的地，若发现他是前往宫廷，则要确认上述所包含的内容（即他在宫廷的目的），以免他可能是应召而来，以便允许他前行。然而，若如贵圣德前述，他前往宫廷是为了请愿和求职，则既不得签署他的信函，也不得收纳他进入共融。\par

众人都说这适当且正确，此规定为他们所认可。\par

此为拉丁文本之第十一教规。\par

% structure: p0140 source=h3 target=subsection reason=relative-heading-rank; base=h2

\subsection{第十二教规（拉丁文本）}

\Person{奥西乌}主教说：但是，亲爱的弟兄们，此处需要一些酌情处理，以防有人来到那些位于大道上的城市，但尚不知晓会议所定之法令。因此，该城市的主教应当告诫他[如此到达的主教]，并指令他从该地派遣他的执事。然而，在此告诫之后，他本人必须返回自己的教区。\par

\Person{奥西乌}此提议在罗马法典中是作为前一教规的附录。希腊文本则完全省略，很可能是因为这似乎是\Person{奥西乌}个人的提议而非主教会议教规，因无记录显示主教会议采纳；或者因为，即使它是法令，也只是暂时性的，即直到教规被充分颁布之前，因此有些人可因不知情而免受所威胁的惩罚。\par

\SourceRecord{Translated by Henry Percival.；Nicene and Post-Nicene Fathers, Second Series；Vol. 14.；Edited by Philip Schaff and Henry Wace.；Buffalo, NY: Christian Literature Publishing Co.,；1900.；<http://www.newadvent.org/fathers/3815.htm>.}
